% Ejemplo completo: muestra todos los elementos del formato.
% Compilar:  pdflatex evaluacion-ejemplo.tex  (dos veces, para "Página X de Y")
\documentclass[carta,doscolumnas,compacta]{evaluacion}
% Opciones: carta|a4 · doscolumnas|unacolumna · normal|compacta|muycompacta
%           solucion (marca la correcta) · sinautoevaluacion

% Colegio, logo, color y docente: de curriculo/config/colegio.yaml y docente.yaml.
% Para otro valor en este documento: \colegio{…} \logo{archivo.png} \colorcolegio{166534} \docente{…}

\asignatura{Pensamiento lógico y tecnología}
\grado{Noveno}
\periodo{Tercer periodo}
\anio{2026}
\docente{Nombre del docente}
% \fecha{2026-11-12}   % vacío = línea para escribir
% \estudiante{...}     % vacío = línea para escribir
\instrucciones{Lee con atención cada pregunta. En las de selección múltiple marca una sola
respuesta (A, B, C o D). Responde las preguntas abiertas con letra clara en las líneas.}

\begin{document}
\encabezado

\begin{preguntas}

\begin{contexto}[Lectura: el reciclaje en el colegio]
\grupo{Responde las preguntas 1 y 2 con base en la siguiente información.}
En una jornada de reciclaje, cada grado recogió botellas plásticas durante una semana.
\begin{tabla}{|l|c|c|}
\encabezados \textbf{Grado} & \textbf{Botellas} & \textbf{Kilogramos} \fila
Séptimo & 120 & 3,0 \fila
Octavo  & 95  & 2,4 \fila
Noveno  & 150 & 3,8 \fila
\end{tabla}
\end{contexto}

\begin{pregunta}[C]
¿Qué grado recogió más botellas?
\opciones{Séptimo}{Octavo}{Noveno}{Todos recogieron lo mismo}
\end{pregunta}

\begin{pregunta}[B]
¿Cuántas botellas recogieron entre los tres grados?
\opciones{355}{365}{375}{345}
\end{pregunta}

\begin{pregunta}[A]
Observa el siguiente programa en Python:
\begin{codigo}[language=Python]
total = 0
for i in range(1, 4):
    total += i * 2   # se suma el doble
print(total)
\end{codigo}
¿Qué valor imprime?
\opciones{\cod{12}}{\cod{6}}{\cod{8}}{\cod{24}}
\end{pregunta}

\begin{pregunta}[D]
¿Cuál de estas afirmaciones sobre los algoritmos es correcta?
\begin{opcioneslista}
\item Un algoritmo solo puede escribirse en un lenguaje de programación.
\item Un algoritmo puede tener infinitos pasos si el problema es difícil.
\item El orden de los pasos de un algoritmo no altera el resultado.
\item Un algoritmo es una secuencia finita y ordenada de pasos que resuelve un problema.
\end{opcioneslista}
\end{pregunta}

\begin{pregunta}[B]
Analiza el siguiente silogismo:
\silogismo{Todo mamífero es vertebrado.}{Todo perro es mamífero.}{Todo perro es vertebrado.}
Sea S = perros, M = mamíferos, P = vertebrados. ¿Qué diagrama representa las premisas?
\opcionesdiagrama
  {\venn[sombreado={SP,MP},escala=0.45]{S,M,P}}
  {\venn[sombreado={S,SP,M,SM},escala=0.45]{S,M,P}}
  {\venn[marcas={SMP},escala=0.45]{S,M,P}}
  {\venn[sombreado={SMP},escala=0.45]{S,M,P}}
\end{pregunta}

\begin{pregunta}[C]
En el diagrama, la región sombreada representa:
\venn[sombreado={AB},titulo={Estudiantes que practican deportes},escala=0.6]{A,B}
\opcioneslargas{Los que practican solo A.}{Los que no practican ninguno.}%
  {Los que practican A y B.}{Los que practican solo B.}
\end{pregunta}

\begin{pregunta}[A]
Según el diagrama, ¿qué se puede concluir?
\venn[sombreado={SM,SMP},marcas={SP},escala=0.6]{S,M,P}
\opciones{Ningún S es M}{Todo S es M}{Algún S es M}{Ningún S es P}
\end{pregunta}

\begin{pregunta}[B]
La gráfica muestra la temperatura de una ciudad durante el día. ¿A qué hora fue máxima?
\begin{center}
\begin{tikzpicture}[scale=0.55]
  \draw[->] (0,0) -- (6.4,0) node[right,font=\tiny] {hora};
  \draw[->] (0,0) -- (0,3.4) node[above,font=\tiny] {°C};
  \draw[thick,acento] (0,0.8) -- (1,1.5) -- (2,2.6) -- (3,3.0) -- (4,2.2) -- (5,1.4) -- (6,1.0);
  \foreach \x/\h in {0/6,1/8,2/10,3/12,4/14,5/16,6/18}
    \draw (\x,0.08) -- (\x,-0.08) node[below,font=\tiny] {\h};
\end{tikzpicture}
\end{center}
\opciones{10:00}{12:00}{14:00}{16:00}
\end{pregunta}

\begin{pregunta}[abierta]
Escribe con tus palabras qué es una variable en programación y da un ejemplo.
\abierta[4]
\modelo{Una variable es un espacio de memoria con un nombre que guarda un valor que puede
cambiar durante la ejecución del programa. Ejemplo: \cod{edad = 15}.}
\end{pregunta}

\begin{pregunta}[abierta]
Observa la imagen y describe el proceso que representa.
\imagen[0.7]{ciclo-del-agua.png}
\abierta[3]
\modelo{El ciclo del agua: evaporación, condensación, precipitación y escorrentía.}
\end{pregunta}

\end{preguntas}

\autoevaluacion
\imprimirclave
\end{document}
